\chapter{Annexes}

%--------------------------------------------------
\section*{Annexe A~: Extrait du Code Python --- Prétraitement et Entraînement}
\addcontentsline{toc}{section}{Annexe A~: Extrait du Code Python}

\begin{small}
\begin{verbatim}
import pandas as pd
import numpy as np
from sklearn.model_selection import train_test_split, cross_val_score
from sklearn.preprocessing import LabelEncoder, StandardScaler
from sklearn.linear_model import LinearRegression, Ridge
from sklearn.ensemble import RandomForestRegressor
from xgboost import XGBRegressor
from sklearn.metrics import mean_absolute_error, mean_squared_error, r2_score

# -- 1. Chargement des données ------------------------------------------------
df = pd.read_csv("yield_df.csv")
print(df.shape)          # (28242, 7)
print(df.isnull().sum()) # Audit des valeurs manquantes

# -- 2. Encodage des variables catégorielles ------------------------------------
le_area = LabelEncoder()
le_item = LabelEncoder()
df['Area']  = le_area.fit_transform(df['Area'])
df['Item']  = le_item.fit_transform(df['Item'])

# -- 3. Séparation features / cible --------------------------------------------
X = df[['Area','Item','Year',
        'average_rain_fall_mm_per_year',
        'pesticides_tonnes','avg_temp']]
y = df['hg/ha_yield']

# -- 4. Split entraînement / test (80/20) --------------------------------------
X_train, X_test, y_train, y_test = train_test_split(
    X, y, test_size=0.2, random_state=42)

# -- 5. Normalisation ----------------------------------------------------------
scaler = StandardScaler()
X_train_sc = scaler.fit_transform(X_train)
X_test_sc  = scaler.transform(X_test)

# -- 6. Entraînement des modèles ----------------------------------------------
models = {
    'Linear Regression' : LinearRegression(),
    'Ridge'             : Ridge(alpha=1.0),
    'Random Forest'     : RandomForestRegressor(
                              n_estimators=200, max_depth=15,
                              random_state=42, n_jobs=-1),
    'XGBoost'           : XGBRegressor(
                              n_estimators=300, max_depth=6,
                              learning_rate=0.05, subsample=0.8,
                              random_state=42),
}

results = {}
for name, model in models.items():
    model.fit(X_train_sc, y_train)
    y_pred = model.predict(X_test_sc)
    results[name] = {
        'MAE' : mean_absolute_error(y_test, y_pred),
        'RMSE': np.sqrt(mean_squared_error(y_test, y_pred)),
        'R2'  : r2_score(y_test, y_pred)
    }
    print(f"{name:20s}  R2={results[name]['R2']:.3f}  "
          f"RMSE={results[name]['RMSE']:.0f}")

# -- 7. Importance des variables (XGBoost) ------------------------------------
xgb = models['XGBoost']
feat_imp = pd.Series(xgb.feature_importances_, index=X.columns)
print(feat_imp.sort_values(ascending=False))
\end{verbatim}
\end{small}

%--------------------------------------------------
\section*{Annexe B~: Statistiques Descriptives du Jeu de Données}
\addcontentsline{toc}{section}{Annexe B~: Statistiques descriptives}

\begin{table}[H]
\centering
\caption{Statistiques descriptives des variables numériques --- Jeu de données AgriYield}
\label{tab:stats-desc}
\renewcommand{\arraystretch}{1.4}
\begin{tabular}{|l|r|r|r|r|r|}
\hline
\rowcolor{maingreen!15}\textbf{Variable} & \textbf{Min} & \textbf{Max} & \textbf{Moyenne} & \textbf{Méd.} & \textbf{Éc.-type} \\
\hline
hg/ha\_yield & 25 & 501\,412 & 28\,764 & 19\,990 & 32\,418 \\
\hline
rain\_fall (mm) & 51 & 3\,240 & 1\,001 & 987 & 588 \\
\hline
pesticides (t) & 0.01 & 367\,778 & 18\,924 & 2\,569 & 48\,341 \\
\hline
avg\_temp ($^\circ$C) & $-$1.4 & 29.7 & 18.2 & 19.7 & 7.1 \\
\hline
Year & 1990 & 2013 & 2001 & 2001 & 6.9 \\
\hline
\end{tabular}
\end{table}

%--------------------------------------------------
\section*{Annexe C~: Hyperparamètres des Modèles Retenus}
\addcontentsline{toc}{section}{Annexe C~: Hyperparamètres des modèles}

\begin{table}[H]
\centering
\caption{Hyperparamètres des modèles entraînés}
\label{tab:hyperparams}
\renewcommand{\arraystretch}{1.4}
\begin{tabularx}{\textwidth}{|>{\bfseries}p{3cm}|X|}
\hline
\rowcolor{maingreen!15}\textbf{Modèle} & \textbf{Hyperparamètres} \\
\hline
Random Forest & \texttt{n\_estimators=200, max\_depth=15, min\_samples\_split=5, random\_state=42} \\
\hline
XGBoost & \texttt{n\_estimators=300, max\_depth=6, learning\_rate=0.05, subsample=0.8, colsample\_bytree=0.8} \\
\hline
MLP & \texttt{hidden\_layers=(128,64,32), activation=ReLU, optimizer=Adam, lr=0.001, epochs=100, batch=64} \\
\hline
Ridge & \texttt{alpha=1.0} \\
\hline
\end{tabularx}
\end{table}

%--------------------------------------------------
\section*{Annexe D~: Architecture de la Plateforme AgriYield}
\addcontentsline{toc}{section}{Annexe D~: Architecture de la plateforme}

\begin{table}[H]
\centering
\caption{Stack technique de la plateforme AgriYield}
\label{tab:stack-tech}
\renewcommand{\arraystretch}{1.4}
\begin{tabularx}{\textwidth}{|>{\bfseries}p{3cm}|p{3cm}|X|}
\hline
\rowcolor{maingreen!15}\textbf{Couche} & \textbf{Technologie} & \textbf{Rôle} \\
\hline
Frontend & React.js / Tailwind & Interface utilisateur, visualisations, upload CSV \\
\hline
Backend API & FastAPI (Python) & Exposition des endpoints de prédiction et simulation \\
\hline
Modèles ML & Scikit-learn, XGBoost & Moteur de prédiction sérialisé (pickle/joblib) \\
\hline
Stockage & PostgreSQL & Historique des prédictions, données utilisateurs \\
\hline
Déploiement & Docker / Cloud & Conteneurisation et scalabilité \\
\hline
\end{tabularx}
\end{table}
